
\section{Introduction}

Secure Enclaves or Trusted Execution Environments (TEEs) are powerful trusted-hardware primitives that can be seen as processes with strengthened privacy and integrity properties.
The security of secure enclaves is often defined in association with a strong threat model where most of the host's computing stack is not trusted.
Not only the host's hypervisor, but also the operating system (OS) and any other piece of unprivileged software collocated on the machine, is considered malicious. 
For simplicity, in the rest of this paper, we will conservatively refer to any untrusted piece of software as the OS.
The security guarantees upheld in the presence of these adversaries are usually two-fold.
First, computation privacy is enforced by different isolation mechanisms, making sure no adversary can extract secrets from the secure enclave.
Second, computation integrity is enforced by an attestation mechanism generating a proof that can be checked by a remote client, attesting that the right application is running on the private data.
Academic TEE proposals date back twenty years to XOM \cite{xom} and AEGIS \cite{aegis} and industry implementations include ARM TrustZone \cite{trustzone04, pinto2019demystifying}, Intel SGX \cite{mckeen2013sgx, costan2016intel} and AMD SEV \cite{sev2020strengthening}.

\begin{figure}[t]
\centering
\includegraphics[width=\columnwidth]{figures/Stack.png}
\caption{Citadel's Computing Stack}
\label{fig:computing_stack}
\end{figure}

Unfortunately, enclaves have not yet become the keystone of secure systems as we have only seen limited instances of real-world deployment.
Several explanations are possible here. 
First, it is challenging to write secure enclave programs, and deployed enclave platforms are still quite recent. 
Another reason could be the many doubts regarding their effective security guarantees.
Indeed, Intel SGX and other platforms have been shown to be vulnerable to microarchitectural side channels and transient execution attacks \cite{chen2019sgxpectre, koruyeh2018spectre, ryan2019hardware, van2020lvi}; in particular, attacks such as those introduced in \cite{van2018foreshadow, van2020sgaxe} were able to not only extract sensitive data from Intel SGX enclaves, but also managed to break the security of the remote attestation mechanism.

These concerns have shifted the focus of trusted computing away from interactive enclaves running critical programs toward scenarios with more modest security goals.
For instance, recent industry technologies like Intel TDX or AMD SEV are designed to isolate entire virtual machines from one another
by partitioning server-class machines and only providing some form of secure boot.
Similarly, Intel recently announced that SGX was deprecated on personal computers and would only be supported on server-class machines.
This setting greatly broadens the TEE attack surface by including entire VMs, hence millions of lines of codes, in the trusted code base, and largely ignores more subtle threats (e.g., no industry group has proposed a platform that would be secure against transient execution attacks).

On the contrary, we believe that enclaves are a promising primitive for highly interactive applications that perform critical security operations, and we should aim at making them secure under more challenging threat models.
Direct applications include executing simple cryptographic operations (like signing, hashing, and encrypting).
Truly secure enclaves available on commodity hardware could be used as a security primitive to enable complex cryptographic protocols such as multiparty computation, federated learning, and confidential computing.
Additionally, this would be possible without paying the cost of heavy cryptography (e.g., Fully Homomorphic Encryption), whose overheads are still multiple orders of magnitude over non-private computation.
These uses can already be found in the few current deployed enclaves applications \cite{signalcontact, signalpin, signalvaluerecovery, zhang2016town, appleplatformsecurity, raj2016ftpm}.

Unfortunately, securing interactive enclaves under a strong threat model that includes transient execution attacks is a challenge that has been mostly left unaddressed. 
Transient execution attacks \cite{Kocher2018spectre, Lipp2018meltdown, behnia2021speculative}
exploit performance optimizations present in modern processors (such as branch predictors or speculative loads) to speculatively access some secrets and change some microarchitectural state in a way that depends on the secret.
A side channel is then used to extract the microarchitectural state and leak the secret.
These attacks are especially hard to defend against as there exist numerous speculation mechanisms and microarchitectural side channels that can be combined in different ways to mount new attacks.

In the specific case of enclaves, most academic enclave platforms \cite{lee2020keystone, bahmani2021cure, nasahl2021hector, costan2016sanctum} circumvent the issue by using in-order processors or physically-isolated cores.
Thus far, only one academic work, MI6 \cite{bourgeat2019mi6} proposes mechanisms to enforce strong enclave isolation on speculative machines, in the presence of transient execution attacks.
MI6 exploits concepts introduced in Sanctum \cite{costan2016sanctum}, a RISC-V enclave platform, and adapts its isolation principles to an out-of-order machine.

MI6 describes and evaluate a large architecture for server-class machines targeted to run enclaves containing non-interactive batch computations. 
Common mechanisms used for efficient data-exchange with an enclave, such as shared memory, were identified as insecure in the MI6 threat model, as they make enclaves vulnerable to Spectre-type attacks.
This effectively prevents any efficient enclave communication operation, like interacting with the parent application to access the network, receiving input, outputting results or executing system calls.
%Effectively, only non-interactive batch computations could be run on this platform.

MI6 considered a 16-core machine equipped with a 16MB last level cache (LLC).
We show in Section \ref{sec:eval} that adapting the large design to a smaller 2-core machine causes critical performance loss, due to static partitioning of the LLC.
Moreover, MI6 enforced conservative resource isolation policies that did not affect the performance of batch-style computations, but hurt more interactive enclaves, preventing, for instance, the use of frequent interrupts.

\textbf{This Work:} In this work, we present a realistic enclave platform, targeting interactive and security-sensitive programs.
First, we solve a key requirement for practical enclaves: we propose and implement the first mechanism for secured shared memory that protects the enclaves from transient execution attacks.
To motivate our proposal, we demonstrate how a naive implementation can be exploited to exfiltrate enclave secrets by mounting a Spectre-style attack on our FPGA insecure baseline.
We also explain that other microarchitectural mechanisms such as noncacheable memory accesses cannot be repurposed to defend these attacks as they do not block speculative address translation.

To thwart these attacks, we modify our processor pipeline to \textit{prevent} enclaves from speculatively accessing \textit{shared} memory, while allowing enclaves to speculatively access enclave memory.
Our approach ensures that no shared microarchitectural state can be modified speculatively depending on an enclave secret.
Hence, any change in the shared microarchitectural state has to be explicitly mandated by the enclave program.
As a result, our platform presents strong microarchitectural isolation guarantees, similar to MI6, while significantly improving enclave expressivity by allowing efficient interactions between an enclave and the outside world. 
We provide a comprehensive security analysis and evaluation for our new secure shared memory mechanism, showing that we considerably improve performance of data exchange.

Second, we provide new hardware-software mechanisms with finer-grained and dynamic isolation mechanisms to improve the full-system performance. 
We propose a new dynamic set partitioning scheme for the Last Level Cache that enables up to 64 security domains in our 1MB LLC, giving us critical performance improvements over static partitioning.
Further, we identify that some context switches do not require  microarchitectural flush, which allows us to enable useful system mechanisms such as frequent time or network interrupts at a reasonable cost.
We provide a security analysis and a performance evaluation of our new mechanisms.

Finally, we build everything presented in the paper as part of our end-to-end FPGA-functional prototype: Citadel. 
Our modified out-of-order 2-core processor can be synthesized and run on an AWS F1 FPGA, boots (untrusted) Linux and makes it possible for user-level applications to launch and safely communicate with secure enclaves executing real applications such as a lightweight Python runtime or a cryptographic library (see Figure \ref{fig:computing_stack}).
We also implement an end-to-end attestation mechanism, making it possible for an enclave to identify itself independently of where the OS placed it in physical memory.
We implement and make all the pieces of software and modified hardware available, such as our secure bootloader, our security monitor (equivalent to processor microcode) and a Linux kernel module to interact with the security monitor.
We also share the implemented hardware isolation mechanisms, like MSHR, LLC arbiter partitioning, and LLC partitioning.
This makes our platform a unique environment for research in the domain of microarchitectural side channels, transient execution attacks and secure enclaves.
We hope researchers will take ownership of the platform and modify it to prototype and test complex hardware-software mechanisms and research ideas.

To summarize, the main contributions of the paper are:
\begin{itemize}
\itemsep-0.1em
\item We propose the first mechanism for enclave shared memory, secure against microarchitectural side channels and transient execution attacks.
\item We refine MI6 isolation mechanisms using hardware-software co-design to improve performance and make the platform practical.
\item To our knowledge, we show the first end-to-end enclave platform running on an FPGA implementation of a speculative processor.
\item We open-source our hardware and software platform at \url{https://github.com/mit-enclaves/citadel}
\end{itemize}